% GENERATED FILE. Edit canonical lessons, then run npm run generate:books.
% canonical-source-hash: fnv1a64:1bcb4eccacb4f4b5
% canonical-lessons: GE-C474-staunen, GE-C474-erstaunt, GE-C474-verwirrt, GE-C474-neugier, GE-C474-sprachlos

\chapter{Staunen und Neugier}
\label{ch:ge-staunen-und-neugier}
\hlchaptermodality{474}

\begin{chapteropening}
\textbf{By the end of this chapter:} \emph{I can talk about surprise, confusion and curiosity.}

Complete the last lesson of chapter 474: I can talk about surprise, confusion and curiosity.
\end{chapteropening}

\section[staunen]{staunen — to be amazed, to marvel}

\label{lesson:GE-C474-staunen}



\begin{quote}
Before the new one: say the German for silly, then the German for to tickle.
\end{quote}

\subsection*{You'll want to know: staunen}
\textbf{staunen} — \textquotedblleft{}to be amazed, to marvel\textquotedblright{}.

\begin{quote}
Die Touristen staunten über die hohen Berge.
\end{quote}

\begin{quote}
\emph{The tourists marvelled at the high mountains.}
\end{quote}

\emph{Staunen über} + accusative — to be amazed at.

\begin{cousinweb}
\emph{staunen} came into standard German from Swiss German \emph{stūnen} \textquotedblleft{}to stiffen, to freeze up\textquotedblright{} (EtymDB): someone who is amazed stands still and stares.
\end{cousinweb}

\begin{grammarlens}[title={What you've built}]
One more word for telling, explaining and describing at B1.
\end{grammarlens}

\subsection*{Your turn}
\begin{itemize}
\raggedright
  \item \emph{Say it:} \emph{staunen}
  \item \emph{Say it:} \emph{staunen}, once more
  \item \emph{Say it:} a sentence of your own with \emph{staunen}
  \item \emph{Recall:} say the German for silly, then the German for to tickle, then say \emph{staunen} again
\end{itemize}

\begin{tcolorbox}[breakable,colback=teal!4,colframe=teal!35!black,title={Before you move on}]
What does staunen mean? (\textquotedblleft{}To be amazed, to marvel\textquotedblright{}.) Say it once more.
\end{tcolorbox}
\section[erstaunt]{erstaunt — astonished}

\label{lesson:GE-C474-erstaunt}



\begin{quote}
Before the new one: say the German for to tickle, then the German for to be amazed.
\end{quote}

\subsection*{You'll want to know: erstaunt}
\textbf{erstaunt} — \textquotedblleft{}astonished\textquotedblright{}.

\begin{quote}
Sie war erstaunt, wie gut ihr Sohn kochen kann.
\end{quote}

\begin{quote}
\emph{She was astonished at how well her son can cook.}
\end{quote}

The participle of \emph{erstaunen} — to astonish.

\begin{cousinweb}
\emph{erstaunt} is built on \emph{er-} and \emph{staunen} \textquotedblleft{}to be amazed\textquotedblright{}, which came from Swiss German \emph{stūnen} \textquotedblleft{}to stiffen\textquotedblright{} (EtymDB).
\end{cousinweb}

\begin{grammarlens}[title={What you've built}]
One more word for telling, explaining and describing at B1.
\end{grammarlens}

\subsection*{Your turn}
\begin{itemize}
\raggedright
  \item \emph{Say it:} \emph{erstaunt}
  \item \emph{Say it:} \emph{erstaunt}, once more
  \item \emph{Say it:} a sentence of your own with \emph{erstaunt}
  \item \emph{Recall:} say the German for to tickle, then the German for to be amazed, then say \emph{erstaunt} again
\end{itemize}

\begin{tcolorbox}[breakable,colback=teal!4,colframe=teal!35!black,title={Before you move on}]
What does erstaunt mean? (\textquotedblleft{}Astonished\textquotedblright{}.) Say it once more.
\end{tcolorbox}
\section[verwirrt]{verwirrt — confused}

\label{lesson:GE-C474-verwirrt}



\begin{quote}
Before the new one: say the German for to be amazed, then the German for astonished.
\end{quote}

\subsection*{You'll want to know: verwirrt}
\textbf{verwirrt} — \textquotedblleft{}confused\textquotedblright{}.

\begin{quote}
Nach der langen Reise war ich müde und verwirrt.
\end{quote}

\begin{quote}
\emph{After the long journey I was tired and confused.}
\end{quote}

The participle of \emph{verwirren} — to confuse.

\begin{cousinweb}
\emph{verwirrt} is built on \emph{ver-} and \emph{wirren} \textquotedblleft{}to move in confusion\textquotedblright{}, from Middle High German \emph{werren} and Old High German \emph{werran}, from Proto-Germanic \emph{werraną} \textquotedblleft{}to confuse\textquotedblright{} (EtymDB).
\end{cousinweb}

\begin{grammarlens}[title={What you've built}]
One more word for telling, explaining and describing at B1.
\end{grammarlens}

\subsection*{Your turn}
\begin{itemize}
\raggedright
  \item \emph{Say it:} \emph{verwirrt}
  \item \emph{Say it:} \emph{verwirrt}, once more
  \item \emph{Say it:} a sentence of your own with \emph{verwirrt}
  \item \emph{Recall:} say the German for to be amazed, then the German for astonished, then say \emph{verwirrt} again
\end{itemize}

\begin{tcolorbox}[breakable,colback=teal!4,colframe=teal!35!black,title={Before you move on}]
What does verwirrt mean? (\textquotedblleft{}Confused\textquotedblright{}.) Say it once more.
\end{tcolorbox}
\section[die Neugier]{die Neugier — curiosity}

\label{lesson:GE-C474-neugier}



\begin{quote}
Before the new one: say the German for astonished, then the German for confused.
\end{quote}

\subsection*{You'll want to know: die Neugier}
\textbf{die Neugier} — \textquotedblleft{}curiosity\textquotedblright{}.

\begin{quote}
Aus Neugier öffnete er das Paket.
\end{quote}

\begin{quote}
\emph{Out of curiosity, he opened the parcel.}
\end{quote}

No plural. Also \emph{die Neugierde}. The adjective is \emph{neugierig} — curious.

\begin{cousinweb}
\emph{Neugier} is built on \emph{neu} \textquotedblleft{}new\textquotedblright{}, inherited from Old High German \emph{niuwi}, from Proto-Germanic \emph{niwjaz} (EtymDB), and \emph{Gier} \textquotedblleft{}greed, craving\textquotedblright{}: a craving for new things.
\end{cousinweb}

\begin{grammarlens}[title={What you've built}]
One more word for telling, explaining and describing at B1.
\end{grammarlens}

\subsection*{Your turn}
\begin{itemize}
\raggedright
  \item \emph{Say it:} \emph{die Neugier}
  \item \emph{Say it:} \emph{die Neugier}, once more
  \item \emph{Say it:} a sentence of your own with \emph{die Neugier}
  \item \emph{Recall:} say the German for astonished, then the German for confused, then say \emph{die Neugier} again
\end{itemize}

\begin{tcolorbox}[breakable,colback=teal!4,colframe=teal!35!black,title={Before you move on}]
What does die Neugier mean? (\textquotedblleft{}Curiosity\textquotedblright{}.) Say it once more.
\end{tcolorbox}
\section[sprachlos]{sprachlos — speechless}

\label{lesson:GE-C474-sprachlos}



\begin{quote}
Before the new one: say the German for confused, then the German for curiosity.
\end{quote}

\subsection*{You'll want to know: sprachlos}
\textbf{sprachlos} — \textquotedblleft{}speechless\textquotedblright{}.

\begin{quote}
Über das teure Geschenk war sie sprachlos.
\end{quote}

\begin{quote}
\emph{She was speechless at the expensive present.}
\end{quote}

\emph{Ich bin sprachlos!} — I'm lost for words!

\begin{cousinweb}
\emph{sprachlos} is built on \emph{Sprache} \textquotedblleft{}speech\textquotedblright{}, inherited from Old High German \emph{sprāhha}, from Proto-Germanic \emph{sprēkijō} (EtymDB), with \emph{-los} \textquotedblleft{}without\textquotedblright{}.
\end{cousinweb}

\begin{grammarlens}[title={What you've built}]
One more word for telling, explaining and describing at B1.
\end{grammarlens}

\subsection*{Your turn}
\begin{itemize}
\raggedright
  \item \emph{Say it:} \emph{sprachlos}
  \item \emph{Say it:} \emph{sprachlos}, once more
  \item \emph{Say it:} a sentence of your own with \emph{sprachlos}
  \item \emph{Recall:} say the German for confused, then the German for curiosity, then say \emph{sprachlos} again
\end{itemize}

\begin{tcolorbox}[breakable,colback=teal!4,colframe=teal!35!black,title={Before you move on}]
What does sprachlos mean? (\textquotedblleft{}Speechless\textquotedblright{}.) Say it once more.
\end{tcolorbox}
